\documentclass[10pt,a4paper]{amsart}
\usepackage[margin=19mm]{geometry}
\usepackage{amsmath,amssymb,mathtools}
\usepackage[scr=rsfs,cal=euler]{mathalfa}
\usepackage{xfrac}
\usepackage[unicode,hidelinks,bookmarks=false]{hyperref}
\newcommand{\ie}{i.e.\ }
\newcommand{\cf}{cf.\ }
\newcommand{\resp}{respectively\ }
\newcommand{\Subsection}{\subsection}
\providecommand{\nobreakdash}{\nobreak\hbox{-}}
\setlength{\emergencystretch}{2em}
\multlinegap0pt
\usepackage{bm}
\usepackage{bbm}
\usepackage{dsfont}
\usepackage{xspace}
\usepackage{cancel}
\usepackage{mathtools}
\usepackage{amsthm}
\makeatletter
\@ifundefined{theorem}{\newtheorem{theorem}{Theorem}}{}
\@ifundefined{corollary}{\newtheorem{corollary}[theorem]{Corollary}}{}
\makeatother
\newcommand{\wt}{\mathsf{wt}}
\title[Cyclic Equalizability of Words and Its Application to Card-B]{Cyclic Equalizability of Words and Its Application to Card-Based Cryptography}
\author{Kazumasa Shinagawa, Koji Nuida}
\date{Source: arXiv:2507.04916v1 (CC BY 4.0)}
\begin{document}
\maketitle

\noindent\emph{Source and licence.} Cyclic Equalizability of Words and Its Application to Card-Based Cryptography. Version arXiv:2507.04916v1 (CC BY 4.0), distributed under the Creative Commons Attribution 4.0 International licence, as recorded in the arXiv licence metadata for this exact version and shown on the abstract page https://arxiv.org/abs/2507.04916v1. Full rights notice: licenses/arxiv-2507.04916-CC-BY-4.0.txt.\par\smallskip

\noindent\textit{Editorial scope.} Counted unit: the Theorem whose source label is \texttt{thm:criterion\_for\_two\_binary\_words} in the article (source0.tex lines 373--376, source label \texttt{thm:criterion\_for\_two\_binary\_words}), delivered together with the article's own complete original proof of it (source0.tex lines 378--457, one \texttt{proof} environment of 80 lines). No mathematics is written, repaired or re-derived: statement and proof are the article's own text line for line. Required context retained uncounted: cor:equalizability\_is\_invariant\_under\_insertion (lines 357--365). Every definition, display and label of those lines is retained unchanged. Omitted from this package and not redistributed: 1--356, 366--372, 377--377, 458--660. Nothing inside a retained excerpt is omitted or altered: every retained body is delivered exactly as the article's lines are written, so no omission receipt of any kind accompanies this package. Citations: the retained excerpts carry no citation command at all, so no bibliography entry of the article is transcribed and no reference is redistributed. Reference adaptation: none was needed and none was made; every reference inside the delivered unit resolves inside it, so no unresolved reference or citation is produced. Printed slips kept literal: no printed word, symbol, hypothesis or number of the counted statement or of its proof was altered. Layout and class: the article's own class is \texttt{article} with packages including fullpage, graphicx, appendix, amsmath, amssymb, amsthm, bm, fancybox, color, cite; the wrapper is the lane's pinned \texttt{amsart} editorial wrapper at 10pt on A4, which restates the theorem-like environments the retained text needs and adds the article's own macro definitions that the retained text needs (\texttt{\textbackslash wt}), with extra wrapper packages (bm, bbm, dsfont, xspace, cancel, mathtools). The delivered numbering is the unit's own. Assets: no retained span contains a figure environment, a graphics-inclusion command, a PDF-page inclusion, a table or a TikZ picture; no image file is copied into this package, the package declares no asset dependency and no image file is redistributed with it. Licence: version arXiv:2507.04916v1 of this article is distributed under the Creative Commons Attribution 4.0 International licence, as recorded in the arXiv licence metadata for this exact version and shown on the abstract page https://arxiv.org/abs/2507.04916v1. A full rights notice is delivered at licenses/arxiv-2507.04916-CC-BY-4.0.txt. Characters: every retained excerpt is pure ASCII, so the delivered engine, XeLaTeX with its default Latin Modern text font, reports no missing character.
\begin{corollary}
    \label{cor:equalizability_is_invariant_under_insertion}
    Suppose that $(w'_1,\dots,w'_k)\in (\Sigma^{n'})^k$ is obtained from $(w_1,\dots,w_k)\in (\Sigma^{n})^k$ by a $\Delta$-deletion for some subset $\Delta \subseteq \Sigma$. 
    Then the following statements are equivalent.
    \begin{enumerate}
        \item $w_1,\dots,w_k$ are $\Delta$-cyclically equalizable.
        \item $w'_1,\dots,w'_k$ are $\Delta$-cyclically equalizable.
    \end{enumerate}
\end{corollary}

\begin{theorem}
    \label{thm:criterion_for_two_binary_words}
    Two binary words $w_1,w_2 \in \{0,1\}^n$ are cyclically equalizable if and only if $\wt(w_1) = \wt(w_2)$.
\end{theorem}

\begin{proof}
    Since the only if part is trivial, we show the if part. 
    Suppose that $\wt(w_1) = \wt(w_2)$. From Corollary~\ref{cor:equalizability_is_invariant_under_insertion}, we can assume that no more $\{0,1\}$-deletion can be applied to $w_1$ and $w_2$, i.e., each column in the matrix representation of the pair $(w_1,w_2)$ is either $(10)^{\mathrm{T}}$ or $(01)^{\mathrm{T}}$.
    Furthermore, since a cyclic shift of the columns of the matrix does not affect the claim, we can also assume that $w_1$ and $w_2$ are of the following form:
    \[
        \begin{split}
            w_1 &= 1^{\mu_0} 0^{\mu_1} 1^{\mu_2} 0^{\mu_3} \cdots 1^{\mu_{2N-2}} 0^{\mu_{2N-1}};\\
            w_2 &= 0^{\mu_0} 1^{\mu_1} 0^{\mu_2} 1^{\mu_3} \cdots 0^{\mu_{2N-2}} 1^{\mu_{2N-1}},
        \end{split}
    \]
    where each $\mu_j$ is a positive integer and it holds $\sum_{j=0}^{N-1} \mu_{2j} = \sum_{j=0}^{N-1} \mu_{2j+1}$ from the assumption $\wt(w_1) = \wt(w_2)$.
    We assume $N \geq 2$ since $w_1$ and $w_2$ are already cyclically equal if $N = 1$. 

    We say that a pair of binary words $(w'_1, w'_2)$ is \emph{admissible} if $\wt(w'_1) = \wt(w'_2)$ and they are of the following form:
    \[
        \begin{split}
            w'_1 &= 1^{\nu'_0} 1^{\mu'_0} 0^{\nu'_1} 0^{\mu'_1} 1^{\nu'_2} 1^{\mu'_2} 0^{\nu'_3} 0^{\mu'_3} \cdots 0^{\nu'_{2N-1}} 0^{\mu'_{2N-1}}; \\
            w'_2 &= 1^{\nu'_0} 0^{\mu'_0} 0^{\nu'_1} 1^{\mu'_1} 1^{\nu'_2} 0^{\mu'_2} 0^{\nu'_3} 1^{\mu'_3} \cdots  0^{\nu'_{2N-1}} 1^{\mu'_{2N-1}},
        \end{split}
    \]
    for some parameters $\nu'_i \geq 0$ and $\mu'_i \geq 1$ ($0 \leq i \leq 2N-1$). 
    The pair $(w_1,w_2)$ in the claim is admissible with $\nu'_i = 0$ and $\mu'_i = \mu_i$. 
    For an admissible pair $(w'_1, w'_2)$ and $0 \leq j \leq 2N-1$, we say that the pair is \emph{consistent} up to the $j$-th slot if the following holds for all $0 \leq i \leq j$:
    \[
    \nu'_i + \mu'_i = \mu'_{(i+1) \bmod 2N} + \nu'_{(i+2) \bmod 2N}.
    \]
    For example, the following pair $(w'_1, w'_2)$ is consistent up to the 1th slot if the underlined subsequences are equal:
    \[
        \begin{split}
            w'_1 &= \underline{1^{\nu'_0} 1^{\mu'_0} 0^{\nu'_1} 0^{\mu'_1}} 1^{\nu'_2} 1^{\mu'_2} 0^{\nu'_3} 0^{\mu'_3} \cdots 0^{\nu'_{2N-1}} 0^{\mu'_{2N-1}}; \\
            w'_2 &= 1^{\nu'_0} 0^{\mu'_0} 0^{\nu'_1} \underline{1^{\mu'_1} 1^{\nu'_2} 0^{\mu'_2} 0^{\nu'_3}} 1^{\mu'_3} \cdots  0^{\nu'_{2N-1}} 1^{\mu_{2N-1}}.
        \end{split}    
    \]
    Note that when $w'_1$ and $w'_2$ are consistent up to the $(2N-3)$-th slot, they are cyclically equal from the following reason. 
    If they are consistent up to the $(2N-3)$-th slot, they are also consistent up to the $(2N-2)$-th slot due to $\wt(w'_1) = \wt(w'_2)$, and furthermore, they are also consistent up to the $(2N-1)$-th slot since they are of equal length. 
    
    In the following, $\{0,1\}$-insertions for $w_1$ and $w_2$ are performed recursively to construct two cyclically equal binary words. 

    First, perform a $\{0,1\}$-insertion to be consistent up to the 0th slot. If $w_1$ and $w_2$ are already consistent up to the 0th slot, i.e., $\mu_0 = \mu_1$, then nothing needs to be done. So we assume that $\mu_0 \neq \mu_1$. 
    If $\mu_0 > \mu_1$, then they will be consistent up to the 0th slot by the following $1$-insertion:
    \[
        \begin{split}
            \widetilde{w_1} &= 1^{\mu_0} 0^{\mu_1} \underline{1^{\mu_0-\mu_1}} 1^{\mu_2} 0^{\mu_3} \cdots 1^{\mu_{2N-2}} 0^{\mu_{2N-1}}; \\
            \widetilde{w_2} &= 0^{\mu_0} 1^{\mu_1} \underline{1^{\mu_0-\mu_1}} 0^{\mu_2} 1^{\mu_3} \cdots 0^{\mu_{2N-2}} 1^{\mu_{2N-1}}.
        \end{split}
    \]
    If $\mu_0 < \mu_1$, then it is enough to do the following $1$-insertion:
    \[
        \begin{split}
            \widetilde{w_1} &= \underline{1^{\mu_1-\mu_0}} 1^{\mu_0} 0^{\mu_1} 1^{\mu_2} 0^{\mu_3} \cdots 1^{\mu_{2N-2}} 0^{\mu_{2N-1}}; \\
            \widetilde{w_2} &= \underline{1^{\mu_1-\mu_0}} 0^{\mu_0} 1^{\mu_1} 0^{\mu_2} 1^{\mu_3} \cdots 0^{\mu_{2N-2}} 1^{\mu_{2N-1}}.
        \end{split}
    \]
    In either case, the current pair $(\widetilde{w_1}, \widetilde{w_2})$ is admissible and consistent up to the 0th slot.

    Next, for $1 \leq j \leq 2N - 3$, when the current pair $(w'_1, w'_2)$ is admissible and consistent up to the $(j-1)$-th slot, we perform $\{0,1\}$-insertions to be consistent up to the $j$-th slot.
    Let $k_0 := \nu'_j + \mu'_j - (\mu'_{j+1} + \nu'_{j+2})$ and $k := |k_0|$.
    If $w'_1$ and $w'_2$ are already consistent up to the $j$-th slot, i.e., $k_0 = 0$, then nothing needs to be done. So we can assume $k_0 \neq 0$.
    If $k_0 > 0$, perform a $0$-insertion of $0^k$ just before $0^{\nu'_{j+2}}$ when $j$ is odd, and a $1$-insertion of $1^k$ just before $1^{\nu'_{j+2}}$ when $j$ is even. 
    The following is an example when $j$ is odd:
    \[
        \begin{split}
            \widetilde{w_1} &= \cdots 1^{\mu'_{j-1}} 0^{\nu'_j} 0^{\mu'_j} 1^{\nu'_{j+1}} 1^{\mu'_{j+1}} \underline{0^{k}} 0^{\nu'_{j+2}} 0^{\mu'_{j+2}} 1^{\nu'_{j+3}} \cdots; \\
            \widetilde{w_2} &= \cdots 0^{\mu'_{j-1}} 0^{\nu'_j} 1^{\mu'_j} 1^{\nu'_{j+1}} 0^{\mu'_{j+1}} \underline{0^{k}} 0^{\nu'_{j+2}} 1^{\mu'_{j+2}} 1^{\nu'_{j+3}} \cdots.
        \end{split}
    \]
    If $k_0 < 0$, for each $0 \leq i \leq j$ such that $i \equiv j \pmod{2}$, 
    a $0$-insertion of $0^k$ just before $0^{\nu'_i}$ is performed when $j$ is odd and a $1$-insertion of $1^k$ just before $1^{\nu'_i}$ when $j$ is even. 
    The following is an example when $j$ is odd:
    \[
        \begin{split}
            \widetilde{w_1} &= 1^{\nu'_0} 1^{\mu'_0} \underline{0^{k}} 0^{\nu'_1} 0^{\mu'_1} 1^{\nu'_2} 1^{\mu'_2} \underline{0^{k}} 0^{\nu'_3} 0^{\mu'_3}\cdots 1^{\mu'_{j-1}} \underline{0^{k}} 0^{\nu'_j} \cdots; \\
            \widetilde{w_2} &= 1^{\nu'_0} 0^{\mu'_0} \underline{0^{k}} 0^{\nu'_1} 1^{\mu'_1} 1^{\nu'_2} 0^{\mu'_2} \underline{0^{k}} 0^{\nu'_3} 1^{\mu'_3}\cdots 0^{\mu'_{j-1}} \underline{0^{k}} 0^{\nu'_j} \cdots.
        \end{split}
    \]
    In either case, the pair $(\widetilde{w_1}, \widetilde{w_2})$ is admissible and consistent up to the $j$-th slot. 
    
    By performing the above procedure up to $j = 2N-3$, we can make $w_1, w_2$ cyclically equal. 
    This completes the proof of Theorem \ref{thm:criterion_for_two_binary_words}.
\end{proof}

\end{document}
